\title{FlashAttention-3: Fast and Accurate Attention with Asynchrony and Low-precision}

\begin{abstract}
The attention mechanism, fundamental to the widespread Transformer framework, constitutes the primary performance constraint for large language models and tasks involving extended contexts. \textsc{FlashAttention} introduced an efficient strategy to accelerate attention computations on GPUs by reducing memory access. Nonetheless, this approach does not fully capitalize on the latest hardware advancements, with \textsc{FlashAttention-2} utilizing only 35\% of the H100 GPU’s capacity. In this work, we present three novel strategies to further enhance attention performance on Hopper GPUs: leveraging the asynchronous operation of Tensor Cores and TMA to (1) concurrently execute computation and data transfer through warp-specialization, (2) alternate between block-level matrix multiplication and softmax computation, and (3) employ block quantization and incoherent processing made possible by hardware-enabled FP8 low-precision support. Our proposed approach, \textsc{FlashAttention-3}, delivers a 1.5-2.0$\times$ speedup on H100 GPUs, achieving up to 740 TFLOPs/s with FP16 (corresponding to 75\% utilization), and approaching 1.2 PFLOPs/s with FP8. Furthermore, we show that FP8 \textsc{FlashAttention-3} yields 2.6$\times$ lower numerical error compared to a standard FP8 attention baseline.
\end{abstract}

\section{Introduction}
\label{sec:intro}

Within the Transformer architecture~\citep{vaswani2017attention}, the chief source of computational expense lies in the attention mechanism, as the calculation of self-attention scores between queries and keys exhibits a quadratic growth with respect to sequence length. Enabling attention to effectively process longer contexts promises to deliver enhanced functionalities—such as the ability to model and perform reasoning across multiple extended documents~\citep{guo2021longt5,shaham2022scrolls,peng2023yarn} and extensive codebases~\citep{roziere2023code, li2023starcoder}—and to support additional modalities like high-resolution images~\citep{chen2022scaling}, audio data~\citep{gulati2020conformer}, and video~\citep{ho2022video}. It also broadens potential use cases, for example, by facilitating user interaction with extensive historical data~\citep{sun2019bert4rec} or managing agent workflows that span lengthy horizons~\citep{yao2022react}. This scenario has sparked considerable research effort focused on increasing the efficiency of attention mechanisms for longer sequences. Approaches include various approximation techniques~\citep{katharopoulos2020transformers,choromanski2020rethinking,tay2020efficient}, advances in software-level optimizations (\citep{rabe2021self,dao2022flashattention,kwon2023efficient}), and even the investigation of alternative architectural designs~\citep{peng2023rwkv,sun2023retentive,gu2023mamba}.

Our research extends the contributions of \citet{dao2022flashattention}, who designed exact-attention algorithms that strategically incorporate the GPU’s architecture and execution features into their overarching algorithmic structure. In their seminal work, \citep{dao2022flashattention}, Dao et al. proposed \textsc{FlashAttention}, which leveraged a tiling method for parallel attention that prevents redundant accesses to the comparatively slow global memory. This was accomplished by merging all attention-related computations within a unified GPU kernel.

Building upon this foundation, \citet{dao2023flashattention2} introduced \textsc{FlashAttention-2}, which reorganized the algorithm to enable parallelization over the sequence length and conducted the core forward-pass computations across partitioned blocks of the key and value matrices. This enhancement sought to boost both GPU work distribution and resource occupancy. Nevertheless, our findings indicate that, despite these improvements, \textsc{FlashAttention-2} achieves suboptimal GPU utilization rates on modern hardware, particularly when compared to finely tuned matrix multiplication (GEMM) kernels—for instance, attaining only 35\% utilization on the Hopper H100 GPU as opposed to the 80-90\% achieved by GEMM kernels. This gap in performance may partly stem from implementation-related factors, such as a failure to employ Hopper-optimized instructions instead of those designed for the Ampere architecture when leveraging Tensor Cores. Recent developments—such as ThunkerKitten~\citep{spector2024thunder} and cuDNN 9~\citep{cudnn9}—demonstrate that adopting Hopper-specific instructions along with tile-oriented computational approaches not only accelerates attention operations but also streamlines their implementation.

At a deeper level, the algorithm implemented by \textsc{FlashAttention-2} follows a basic synchronous execution model and does not intentionally leverage asynchronous execution or low-precision arithmetic in its structure. Asynchrony stems from specialized hardware components that accelerate key operations in machine learning tasks: for instance, certain hardware blocks are dedicated to matrix multiplication (such as Tensor Cores) or to memory transfer operations (like the Tensor Memory Accelerator -- TMA), functioning independently from the rest of the CUDA cores that handle control flow, integer arithmetic, or general floating-point operations. The adoption of low-precision formats—FP8 in Hopper, FP4 in Blackwell, as well as earlier advances with FP16 (introduced with Pascal in 2017) and BF16 (with Ampere in 2020)—is an established strategy to achieve two to four times higher throughput without increasing power consumption or silicon area. \Cref{subsec:hardware} provides an overview of the hardware features introduced in Hopper relevant to these optimizations.

The main technical obstacle is to reengineer \textsc{FlashAttention-2} in a way that fully exploits these hardware advancements. Employing asynchrony entails designing execution paths where matrix multiplication and softmax operations proceed in parallel, despite their inherent dependency, while adopting low-precision computations demands careful attention to control quantization errors—this is particularly crucial for rare or outlier features in large language models, as discussed by~\citet{dettmers2208llm, sun2024massive}.

In pursuit of this goal, we introduce \textsc{FlashAttention-3}, presenting a combination of three novel techniques designed to enhance efficiency on recent GPU platforms.\footnote{Our findings are framed with respect to NVIDIA's Hopper architecture, though our method can be applied to any GPU that offers adequate asynchronous processing and supports low-precision operations.}

\begin{enumerate}

\item \textbf{Producer-Consumer asynchrony:} We introduce a warp-specialized approach to software pipelining that capitalizes on the asynchronous behavior between data transfer operations and computations on Tensor Cores. By distributing data producers and consumers across distinct warps, our scheme enhances the capacity of the algorithm to effectively mask both memory access and instruction issue delays.

\item \textbf{Hiding softmax under asynchronous block-wise GEMMs:} We achieve concurrency between the relatively less compute-intensive non-GEMM components of softmax, such as floating point fused multiply-adds and exponentials, and the asynchronous WGMMA GEMM instructions. To enable this overlap, we revise the \textsc{FlashAttention-2} framework, eliminating specific ordering constraints between softmax processing and GEMMs. For instance, in our algorithm’s two-stage pipeline, softmax computations are performed on one tile of the scores matrix while, in parallel, WGMMA asynchronously processes the subsequent tile.

\item \textbf{Hardware-accelerated low-precision GEMM:} We revise the forward computation method to utilize FP8 Tensor Cores for GEMM operations, yielding close to a twofold improvement in achieved TFLOPs/s. To accomplish this, we address the discrepancies in memory layout expectations for WGMMA, particularly concerning FP32 accumulators and FP8 operand blocks. Methods such as block quantization and incoherent processing are employed to alleviate the precision loss that accompanies the adoption of FP8 arithmetic.
\end{enumerate}

For empirical validation, we conducted benchmarks of \textsc{FlashAttention-3} on the H100 SXM5 GPU across various configurations. Our results indicate that (1) in the forward pass, FP16 provides a speedup of 1.5-2.0$\times$ over \textsc{FlashAttention-2} (achieving up to 740 TFLOPs/s), and a 1.5-1.75$\times$ improvement in the backward pass, (2) FP8 attains performance near 1.2 PFLOPs/s, and (3) when sequence lengths are large, FP16 delivers superior throughput while FP8 remains competitive\footnote{Specifically, when the head dimension is 64, \textsc{FlashAttention-3} FP8 surpasses alternatives; for head dimensions 128 and 256, it matches leading performance in the absence of causal masking, and slightly lags when causal masking is enabled.} with NVIDIA’s cuDNN library’s top attention implementation. Additionally, we confirm that FP16 in \textsc{FlashAttention-3} incurs equivalent numerical error as \textsc{FlashAttention-2} and surpasses standard attention implementations, as intermediate calculations (such as softmax rescaling) are performed in FP32 precision. Finally, for input distributions featuring outlier features, FP8 \textsc{FlashAttention-3}—using block quantization alongside incoherent processing—delivers a 2.6$\times$ accuracy improvement over standard attention employing per-tensor quantization.

We are releasing \textsc{FlashAttention-3} as open-source under a permissive license\footnote{\textsc{FlashAttention-3}
  is available at \texttt{https://github.com/Dao-AILab/flash-attention}} and intend to incorporate it into the PyTorch and Hugging Face ecosystems, aiming to maximize accessibility for both researchers and developers.

\section{Background: Multi-Head Attention and GPU Characteristics}
\label{sec:background}

\subsection{Multi-Head Attention}
\label{subsec:multi_head_attn}

Consider the query, key, and value matrices $\mathbf{Q}, \mathbf{K}, \mathbf{V} \in \mathbb{R}^{N \times d}$ for a single attention head, with $N$ denoting the sequence length and $d$ representing the dimensionality per head. The attention output $\mathbf{O}$ is then calculated as follows:

\begin{equation*}
  \mathbf{S} = \alpha \mathbf{Q} \mathbf{K}^\top \in \mathbb{R}^{N \times N}, \quad \mathbf{P} = \mathrm{softmax}(\mathbf{S}) \in \mathbb{R}^{N \times N}, \quad \mathbf{O} = \mathbf{P}\mathbf{V} \in \mathbb{R}^{N \times d},
\end{equation*}

Here, $\mathrm{softmax}$ operates on each row independently, and it is common to use $\alpha = 1/\sqrt{d}$ as the scale. To mitigate issues related to numerical instability when applying the exponential, we typically offset $\mathbf{S}$ by subtracting $\mathrm{rowmax}(\mathbf{S})$. In the multi-head attention (MHA) framework, each head independently learns separate query, key, and value projections, with these computations executed concurrently across different heads and batches to yield the complete output tensor.

Consider $\phi$ as a scalar loss function, and denote the gradient operator by $\mathbf{d}(-) = \partial \phi / \partial (-)$. Provided the gradient with respect to the output, $\mathbf{dO} \in \mathbb{R}^{N \times d}$, the gradients with respect to $\mathbf{Q}$, $\mathbf{K}$, and $\mathbf{V}$, namely $\mathbf{dQ}$, $\mathbf{dK}$, and $\mathbf{dV}$, are determined via the chain rule as shown below:

\begin{align*}
  \mathbf{dV} &= \mathbf{P}^\top \mathbf{dO} \in \mathbb{R}^{N \times d} \\
  \mathbf{dP} &= \mathbf{dO} \mathbf{V}^\top \in \mathbb{R}^{N \times N} \\
  \mathbf{dS} &= \mathrm{dsoftmax} (\mathbf{dP}) \in \mathbb{R}^{N \times N} \\
  \mathbf{dQ} &= \alpha \mathbf{dS} \mathbf{K} \in \mathbb{R}^{N \times d} \\
  \mathbf{dK} &= \alpha \mathbf{dS}^\top \mathbf{Q} \in \mathbb{R}^{N \times d},
\end{align*}

In this context, we obtain $\mathbf{d}s = (\mathrm{diag}(p) - p p^\top)\mathbf{d}p$, where $p = \mathrm{softmax}(s)$ is defined with respect to the vector $s$. The notation $\mathrm{dsoftmax}(\mathbf{dP})$ refers to applying this expression to each row individually. This operation also allows for efficient parallelization over the number of heads and batches when performing the backward pass in MHA.

\subsection{GPU hardware characteristics and execution model}
\label{subsec:hardware}

We outline the components of the GPU execution model pertinent to \textsc{FlashAttention-3}, concentrating particularly on the NVIDIA Hopper architecture as a specific realization of these principles.

\paragraph{Memory hierarchy:} The GPU utilizes a hierarchical arrangement of memory types, where increased storage capacity typically corresponds to lower bandwidth, as detailed in \cref{tab:gpu-hierarchy}.\footnote{\citet{luo2024benchmarking} cites a shared memory bandwidth of 128 bytes per clock cycle for each SM; this value is then scaled by 132 SMs and a boost clock frequency of 1830 MHz.} The largest tier, global memory (GMEM) or HBM, is off-chip DRAM that all streaming multiprocessors (SMs) can access. Data residing in GMEM is automatically cached into an on-chip L2 cache for faster access. Additionally, every SM possesses a relatively small, highly banked on-chip cache termed shared memory (SMEM), which the programmer explicitly manages. At the lowest level in the hierarchy is the register file, housed within each SM.

\paragraph{Thread hierarchy:} In the GPU programming paradigm, threads are structured into several nested levels of organization. Starting from the most granular, the hierarchy consists of individual threads, followed by warps containing 32 threads each, then warpgroups comprising 4 sequential warps, progressing to threadblocks (also referred to as cooperative thread arrays or CTAs), threadblock clusters (introduced with Hopper), and ultimately grids at the highest tier.

The two hierarchies exhibit a strong interconnection. Threads belonging to a single CTA are concurrently scheduled on the same SM, while CTAs that reside within one cluster are co-scheduled on the same GPC. All threads inside a CTA can directly access SMEM, while each thread individually has access to up to 256 private registers (RMEM).

\paragraph{Asynchrony and warp-specialization:}

GPUs function as throughput-oriented processors, leveraging parallelism and asynchronous operations to mask both memory and computation delays. For asynchronous memory transfers between GMEM and SMEM, Hopper introduces the Tensor Memory Accelerator (TMA), a specialized hardware component \cite[\S7.29]{cuda}. In addition, Hopper diverges from previous designs like Ampere by equipping its Tensor Core—accessible using the warpgroup-level WGMMA instruction \cite[\S9.7.14]{ptx}—with asynchronous capability, enabling it to read inputs straight from shared memory.

The provision of hardware-level asynchrony enables warp-specialized kernels, assigning each warp within a CTA exclusively to either a producer or consumer role—dedicated to data transfer or computation, respectively. Such specialization broadly enhances the compiler’s capacity to construct highly efficient instruction schedules \citep{warp-specialization-2011}. Furthermore, Hopper introduces the capability to dynamically redistribute registers among warpgroups using \verb|setmaxnreg| \cite[\S9.7.17.1]{ptx}, thereby permitting MMAs-focused warps to access a greater portion of RMEM in contrast to those limited to performing TMA, where only one thread is necessary.

\paragraph{Low-precision number formats:}
\label{sec:low-precision-gpu}
Contemporary GPUs incorporate dedicated hardware modules to boost the performance of low-precision calculations. For instance, on Hopper architectures, the WGMMA instruction leverages FP8 Tensor Cores, achieving twice the per-SM throughput relative to computations carried out in FP16 or BF16 formats.

Nevertheless, properly utilizing FP8 WGMMA requires awareness of specific operand layout requirements. Consider the GEMM operation $A \times B^{\top}$, where $A$ is an $M\times K$ matrix and $B$ is an $N\times K$ matrix. An operand is described as \emph{mn-major} if its storage is contiguous along the outer $M$ or $N$ axis, whereas it is termed \emph{k-major} if contiguity is along the inner $K$-axis. In the context of FP16 WGMMA, operands stored in shared memory (SMEM) can utilize either the mn-major or k-major arrangement. Conversely, FP8 WGMMA exclusively supports the k-major ordering for inputs. Additionally, in applications like attention mechanisms, where chaining multiple GEMMs within a single kernel is advantageous, mismatches between the FP32 accumulator layout and the FP8 operand format can hinder the sequential execution of FP8 WGMMAs.

When considering attention mechanisms, these constraints on layout necessitate specific adjustments to the construction of an FP8 algorithm, as detailed in \cref{sec:algofp8}.

\subsection{Standard Attention and Flash Attention} As described by \citet{dao2022flashattention}, the term \textbf{standard attention} refers to a GPU-based attention computation in which the intermediate matrices $\mathbf{S}$ and $\mathbf{P}$ are stored in HBM. The central innovation of \textsc{FlashAttention} is to introduce a local softmax reduction, which circumvents costly intermediate memory accesses by combining the entire attention operation within a single computational kernel. This localized softmax operation is performed in lines \ref{code-ws:softmax_start}-\ref{code-ws:softmax_end} of the consumer mainloop in \cref{alg:flash3_wgmma_ws_only}, alongside the appropriate rescaling of $\mathbf{O}$'s block components. A concise proof showing that this approach correctly produces $\mathbf{O}$ is detailed in \cite[\S 2.3.1]{dao2023flashattention2}.

\section{FlashAttention-3: Algorithm}
\label{sec:algo}

In this section, we present the \textsc{FlashAttention-3} algorithm. Our discussion will primarily address the forward computation, deferring details of the backward pass to~\cref{sec:algo_ws_bwd}. We begin by illustrating the integration of warp-specialization and a circular SMEM buffer into the foundational approach established by \textsc{FlashAttention-2}. Next, we detail the manner in which WGMMA’s asynchronous capabilities enable the construction of a two-stage pipeline that overlaps GEMM and softmax computations. Lastly, we outline the necessary adaptations for FP8 support, covering both layout compatibility and accuracy considerations through block quantization and incoherent data handling.

\subsection{Producer-Consumer asynchrony through warp-specialization and pingpong scheduling}
\label{sec:algo_ws}

\paragraph{Warp-specialization}
Similar to \textsc{FlashAttention-2}, the forward computation in \textsc{FlashAttention-3} is highly parallelizable across the batch dimension, across heads, and along the length of the query sequence. Therefore, it is sufficient to present the algorithm from a CTA perspective, wherein each CTA handles a specific tile $\mathbf{Q}_i$ of the query matrix and generates the matching output tile $\mathbf{O}_i$. To keep the exposition clear, we start by describing the warp-specialization approach assuming a circular SMEM buffer without incorporating the overlap of GEMM and softmax computations. Let $d$ denote the head size and $N$ denote the total sequence length. Choosing a block size $B_r$ for the queries partitions $\mathbf{Q}$ into $T_r = \lceil \frac{N}{B_r} \rceil$ distinct blocks $\mathbf{Q}_1, .., \mathbf{Q}_{T_r}$.

\begin{algorithm}[H]
    \caption{\small\label{alg:flash3_wgmma_ws_only}\textsc{FlashAttention-3} forward pass \textbf{without} intra-consumer overlapping -- CTA view}
    \begin{algorithmic}[1]
\REQUIRE Matrices $\mathbf{Q}_i \in \mathbb{R}^{B_r \times d}$ along with $\mathbf{K}, \mathbf{V} \in \mathbb{R}^{N \times d}$ residing in HBM, key block size $B_c$, and $T_c = \lceil \frac{N}{B_c} \rceil$.
\STATE Create and configure a pipeline to control barrier synchronization using a $s$-stage circular SMEM buffer.
\IF {current context is a producer warpgroup}
\STATE Free up a specified quantity of registers as pre-arranged.
\STATE Move $\mathbf{Q}_i$ from HBM to shared memory.
\STATE After the transfer, signal to the consumer that $\mathbf{Q}_i$ has been loaded.
\FOR{$0 \le j < T_c$}
    \STATE Await confirmation that the $(j\,\%\,s)$-th buffer stage has been processed.
    \STATE Transfer $\mathbf{K}_j, \mathbf{V}_j$ from HBM to the $(j\,\%\,s)$-th position in shared memory.
    \STATE Once loaded, signal to consumers that $\mathbf{K}_j, \mathbf{V}_j$ are available.
\ENDFOR
\ELSE 
\STATE Allocate registers as determined by the consumer warp count.
\STATE Within on-chip memory, set initial values $\mathbf{O}_i = (0) \in \mathbb{R}^{B_r \times d}$ and initialize $\ell_i, m_i = (0), (-\infty) \in \mathbb{R}^{B_r}$.
\STATE Pause execution until $\mathbf{Q}_i$ becomes available in shared memory.
\FOR{$0 \le j < T_c$}
\STATE Wait until $\mathbf{K}_j$ has been copied to shared memory.
\STATE Perform $\mathbf{S}_i^{(j)} = \mathbf{Q}_i \mathbf{K}_j^T$ (SS-GEMM). Commit and synchronize. \label{alg:ws_only_gemm1}
\STATE Store $m_i^{\mathrm{old}} = m_i$ and update $m_i$ with $\mathrm{max}(m_i^{\mathrm{old}}, \mathrm{rowmax}(\mathbf{S}_i^{(j)}))$. \label{code-ws:softmax_start}
\STATE Calculate $\widetilde{\mathbf{P}}_i^{(j)} = \mathrm{exp}(\mathbf{S}_i^{(j)} - m_i)$, and update $\ell_i = \mathrm{exp}(m_i^{\mathrm{old}} - m_i) \ell_i + \mathrm{rowsum}(\widetilde{\mathbf{P}}_i^{(j)})$. \label{code-ws:softmax_end}
\STATE Wait until $\mathbf{V}_j$ is accessible in shared memory.
\STATE Compute the update $\mathbf{O}_i = \mathrm{diag}(\mathrm{exp}(m_i^{\mathrm{old}} - m_i))^{-1} \mathbf{O}_i + \widetilde{\mathbf{P}}_i^{(j)} \mathbf{V}_j$ (RS-GEMM). Commit and synchronize. \label{alg:ws_only_gemm2}
\STATE Allow the producer to reclaim the $(j\,\%\,s)$-th buffer stage.
\ENDFOR
\STATE Finalize with $\mathbf{O}_i = \mathrm{diag}(\ell_i)^{-1} \mathbf{O}_i$ and $L_i = m_i + \log(\ell_i)$.
\STATE Output $\mathbf{O}_i$ and $L_i$ to HBM to serve as the $i$-th blocks of $\mathbf{O}$ and $L$ respectively.
\ENDIF
\end{algorithmic}
\end{algorithm}

In our deployment of \cref{alg:flash3_wgmma_ws_only} on Hopper, we utilize \verb|setmaxnreg| for managing (de)allocation of registers, employ TMA for transferring data for $\mathbf{Q}_i$ and $\{ \mathbf{K}_j, \mathbf{V}_j \}_{0 \leq j < T_c}$, and leverage WGMMA to perform GEMMs within the main consumer loop. Here, the prefixes SS and RS specify whether the initial operand resides in shared memory or in the register file, respectively. When analyzing the control flow of \cref{alg:flash3_wgmma_ws_only}, it is important to recognize that TMA load instructions operate asynchronously, ensuring that the dispatch of a load does not depend on the completion of other outstanding loads. Additionally, within the main loop of the producer, no wait instructions are issued during the initial $s$ iterations since these serve to prime the buffer.

\paragraph{Pingpong scheduling} The asynchrony provided by WGMMA and TMA, coupled with warp specialization, makes it possible to overlap the softmax operation being performed by one warpgroup with the GEMM performed by another. This approach is especially relevant considering the disparity in throughput between matmul and non-matmul operations on state-of-the-art hardware. For instance, the H100 SXM5 GPU delivers 989 TFLOPS for FP16 matmul, whereas it achieves merely 3.9 TFLOPS for special functions like exponentiation\footnote{According to the CUDA programming guide, each streaming multiprocessor (SM) can execute 16 special function operations per cycle. Given 132 SMs running at a clock speed of 1830 MHz, this results in a throughput of 3.9 TFLOPS for special functions.}. In the attention forward pass with FP16 and a head dimension of 128, the number of FLOPS needed for matmul is 512 times greater than for exponentials, but exponential instructions have 256 times lower throughput. Consequently, evaluating exponentials may account for up to 50\% of total execution time, compared to matmul. The situation becomes even more pronounced when using FP8, where matmul throughput doubles, but throughput for exponential functions remains unchanged.

Because the exponential operation is handled by a dedicated hardware unit known as the multi-function unit, optimal scheduling would overlap the exponential computation with the matrix multiplication on the Tensor Cores. To accomplish this, we employ synchronization barriers (\texttt{bar.sync} instructions) to ensure that the GEMMs (specifically, GEMM1, which computes $\mathbf{P} \mathbf{V}$ in the current iteration, and GEMM0, which computes $\mathbf{Q} \mathbf{K}^\top$ in the subsequent iteration) in warpgroup 1 are executed prior to those in warpgroup 2. This approach allows warpgroup 1 to carry out the softmax computation concurrently with warpgroup 2’s GEMMs. The roles then alternate, such that warpgroup 2 executes the softmax while warpgroup 1 proceeds with the GEMMs—a pattern we refer to as "pingpong" scheduling, visualized in~\cref{fig:pingpong_scheduling}. Although in real-world scenarios the pingpong scheduling pattern is not as perfectly orderly as shown in the figure, our measurements consistently show that this technique boosts throughput (for example, achieving an increase from 570 TFLOPS to 620–640 TFLOPS on FP16 forward passes with a head dimension of 128 and sequence length of 8192).

\paragraph{Attention variants} In the case of multi-query attention~\citep{shazeer2019fast} and grouped query attention~\citep{ainslie2023gqa}, our implementation aligns with the methodology employed by \textsc{FlashAttention-2}, modifying tensor indexing in such a way that $\mathbf{K}$ and $\mathbf{V}$ are not redundantly stored in HBM.

\subsection{Intra-warpgroup overlapping GEMMs and softmax}

It is possible to interleave certain instructions from the softmax operation with those from the GEMMs, even when restricted to a single warpgroup. Here, we outline a specific method for achieving such overlap.

Within the attention algorithm, certain operations in the inner loop exhibit sequential dependencies, which restrict the degree of parallelism achievable within a single iteration. For instance, the computation of the (local) softmax (specified in lines \ref{code-ws:softmax_start} to \ref{code-ws:softmax_end}) is contingent on the output $\mathbf{S}_i^{(j)}$ produced by the initial GEMM, while the subsequent GEMM uses $\widetilde{\mathbf{P}}_i^{(j)}$ as an input, which is the result of the softmax. As a consequence, the presence of wait statements at lines \ref{alg:ws_only_gemm1} and \ref{alg:ws_only_gemm2} in \cref{alg:flash3_wgmma_ws_only} enforces a serial order between the execution of the GEMMs and softmax. Nevertheless, it is possible to alleviate these dependencies by introducing pipelining between iterations, achieved by incorporating additional buffering in registers. Building on this concept, we introduce a two-stage\footnote{It is important to highlight that the stage count for the overlapping approach is capped by, though not necessarily equal to, the parameter $s$, which denotes the number of stages in the circular SMEM buffer.} pipelined GEMM-softmax algorithm:

\begin{algorithm}[H]
    \caption{\small\label{alg:flash3_wgmma}\textsc{FlashAttention-3} Consumer Warpgroup Forward Computation}
    \begin{algorithmic}[1]
      \REQUIRE  Input matrices $\mathbf{Q}_i \in \mathbb{R}^{B_r \times d}$ and $\mathbf{K}, \mathbf{V} \in \mathbb{R}^{N \times d}$ residing in HBM; block size for keys $B_c$, where $T_c = \lceil \frac{N}{B_c} \rceil$.
      \STATE Dynamically allocate a specified set of registers depending on the number of consumer warps present.
      \STATE Within on-chip memory, set initial values: $\mathbf{O}_i = (0) \in \mathbb{R}^{B_r \times d}$, and vectors $\ell_i, m_i = (0), (-\infty) \in \mathbb{R}^{B_r}$.
      \STATE Stall until $\mathbf{Q}_i$ and $\mathbf{K}_0$ are available in shared memory.
      \STATE Employ WGMMA to compute $\mathbf{S}_{\mathrm{cur}} = \mathbf{Q}_i \mathbf{K}_0^T$. Commit operation and pause until it completes.
      \STATE Free the $0$th buffer stage dedicated to $\mathbf{K}$.
      \STATE Update $m_{i}$, $\tilde{\mathbf{P}}_{\mathrm{cur}}$, and $\ell_{i}$ using $\mathbf{S}_{\mathrm{cur}}$, with corresponding rescaling of $\mathbf{O}_i$.
      \FOR{$1 \le j < T_c - 1$}
        \STATE Await readiness of $\mathbf{K}_j$ in shared memory.
        \label{code:mainloop_start}
        \STATE Calculate $\mathbf{S}_{\mathrm{next}} = \mathbf{Q}_i \mathbf{K}_{j}^T$ by invoking WGMMA; commit execution but do not block. \label{code:first_wgmma}
        \STATE Await until $\mathbf{V}_{j-1}$ is accessible in shared memory.
        \STATE Via WGMMA, perform $\mathbf{O}_{i} = \mathbf{O}_{i} + \tilde{\mathbf{P}}_{\mathrm{cur}} \mathbf{V}_{j-1}$; commit this operation without stalling. \label{code:second_wgmma}
        \STATE Suspend until the result of WGMMA for $\mathbf{Q}_i \mathbf{K}_{j}^T$ becomes available.
        \STATE From $\mathbf{S}_{\mathrm{next}}$, recalculate $m_{i}$, $\tilde{\mathbf{P}}_{\mathrm{next}}$, and $\ell_{i}$. \label{code:softmax}
        \STATE Wait for completion of WGMMA corresponding to $\tilde{\mathbf{P}}_{\mathrm{cur}} \mathbf{V}_{j-1}$, then adjust scaling of $\mathbf{O}_i$.
        \STATE Release the buffer slots $(j\,\%\,s)$ for $\mathbf{K}$ and $(j-1\,\%\,s)$ for $\mathbf{V}$, respectively.
        \STATE Transfer $\mathbf{S}_{\mathrm{next}}$ data into $\mathbf{S}_{\mathrm{cur}}$ for the next iteration.
        \label{code:mainloop_end}
      \ENDFOR \STATE Wait until $\mathbf{V}_{T_c - 1}$ has finished loading into shared memory.
      \STATE Execute the update $\mathbf{O}_{i} = \mathbf{O}_{i} + \tilde{\mathbf{P}}_{\mathrm{last}} \mathbf{V}_{T_c - 1}$ utilizing WGMMA, committing and waiting until it finalizes.

\STATE Finalization: Adjust $\mathbf{O}_{i}$ by the factor $m_{i}$. Determine $L_{i}$ utilizing both $m_{i}$ and $\ell_{i}$.
Store $\mathbf{O}_{i}$ and $L_{i}$ in HBM as the $i$-th segments of $\mathbf{O}$ and $L$.

\cref{alg:flash3_wgmma} substitutes the consumer path presented in \cref{alg:flash3_wgmma_ws_only}, thereby forming the full \textsc{FlashAttention-3} algorithm at FP16 precision. In this context, WGMMA serves as shorthand for asynchronous GEMM. Throughout the mainloop (spanning lines \ref{code:mainloop_start} through \ref{code:mainloop_end}), the second WGMMA operation in iteration $j$ (line \ref{code:second_wgmma}) executes concurrently with the softmax computations of the subsequent iteration $j+1$ (line \ref{code:softmax}).

Although the pipelined design depicted above promises notable theoretical speedups, its practical implementation is subject to several important considerations:

\paragraph{Compiler reordering} Although the pseudocode outlines an ideal execution sequence, the NVCC compiler frequently modifies instruction ordering to enhance efficiency. Such reordering can interfere with the intended interleaving of WGMMA and non-WGMMA instructions, thereby reducing or negating the anticipated performance benefits. Examination of the generated SASS code reveals that, in this case, the compiler does produce the expected overlapping execution (see Section~\ref{sec:2-stage-sass}).

\paragraph{Register pressure} Achieving peak performance necessitates avoiding register spilling. However, implementing the 2-stage pipeline increases register requirements to accommodate both intermediate results and necessary context for the overlapping stages. Specifically, a separate $\mathbf{S}_{\mathrm{next}}$ must be held in registers, adding to the register allocation by $B_r \times B_c \times \text{sizeof}(\text{float})$ per threadblock. Such heightened register utilization may compete with optimizations such as increasing block sizes, which also demand additional registers. In real applications, the optimal configuration often involves balancing these competing factors according to profiling insights.

\paragraph{3-stage pipelining} Building upon the previously discussed 2-stage pipeline, a 3-stage variation is proposed to facilitate concurrent execution of the second WGMMA and the softmax computation. This strategy could result in even greater exploitation of Tensor Core resources, but necessitates allocating yet more registers for the expanded pipeline depth. The tension between maximizing tile size and deepening the pipeline thus becomes more pronounced. For further technical detail and experimental analysis of this approach, refer to ~\cref{sec:3-stage}.

\subsection{Low-precision with FP8}
\label{sec:algofp8}

\textbf{Efficiency: layout transformations.} Executing the forward pass of \textsc{FlashAttention-3} using FP8 precision introduces unique challenges related to layout alignment, which are not present when operating in FP16.

To begin, it is important to recognize that the input tensors $\mathbf{Q}$, $\mathbf{K}$, and $\mathbf{V}$ are conventionally stored with the head dimension as the fastest-varying (i.e., contiguous) axis. However, to fulfill the k-major layout requirement imposed by FP8 WGMMA for the second GEMM, it is necessary for $\mathbf{V}$—specifically, for the portions of $\mathbf{V}$ brought into SMEM—to have the sequence length dimension contiguous. Because the TMA load operation cannot modify which dimension is contiguous, we are presented with two alternatives: (1) pre-emptively transpose $\mathbf{V}$ in GMEM prior to computation, or (2) perform a transpose on the loaded $\mathbf{V}$ tiles within the kernel after they have been brought into SMEM. For strategy (1), the transpose could be (1a) integrated into the epilogue of an upstream operation, such as the rotary embedding, or (1b) carried out using a separate pre-processing kernel dedicated to transposition\footnote{A high-performance transpose kernel can reach throughput approaching device bandwidth~\citep{colfax_cutlass_transpose_2024}.}, effectively swapping the strides between the sequence length and head axes. Nonetheless, option (1a) poses integration challenges in generic library contexts, and option (1b) is impractical in memory bandwidth-limited scenarios typical of inference workloads.

For FP8 \textsc{FlashAttention-3}, we select approach (2). Within the kernel, the transpose operation leverages the LDSM (\verb|ldmatrix|) and STSM (\verb|stmatrix|) instructions. These allow a thread warp to cooperatively transfer data between SMEM and RMEM in 128-byte blocks.\footnote{The PTX specification refers to LDSM/STSM as mechanisms for copying $8 \times 8$ matrices of 16-bit values \cite[\S 9.7.13.4.15-16]{ptx}. Nevertheless, it is possible to utilize these instructions for FP8 by packing pairs of 8-bit elements. The transpose versions, however, are unable to individually unpack the paired 8-bit values, which means that additional register operations must be inserted between LDSM and STSM in order to achieve an actual tile-wise transpose; details are elided here.} The LDSM/STSM operations are highly efficient in terms of register usage, enabling their use within the producer warpgroup, and they can convert data layouts in the process of memory copying. Furthermore, after the initial iteration, we are able to schedule the transposition of the subsequent $\mathbf{V}$ tile so that it overlaps with the two WGMMAs computations involving the previous $\mathbf{V}$ and the current $\mathbf{K}$ tile.

Second, we note that, in contrast to FP16, the FP32 accumulator within an FP8 WGMMA possesses a distinct memory layout compared to the expected register-based layout of operand A. Illustrations of sections of these two layouts are provided in \cref{fig:rmem_accum} and \cref{fig:rmem_operand}, which demonstrate the sequence in which each thread stores register entries. Leveraging byte permute instructions enables us to convert the accumulator produced by the first WGMMA into a configuration compatible both with a subsequent WGMMA and with the format of the $\mathbf{V}$ tile output by the in-kernel transpose. Concretely, as shown in \cref{fig:rmem_accum}, we reorder the entries into the pattern
$$\{ \verb|d0 d1 d4 d5 d2 d3 d6 d7| \},$$
applying this specific register shuffle to each contiguous group of 8 bytes. Regarding the structural organization of the $\mathbf{P}$ tile, this operation permutes its columns—for example, columns $0189$ are mapped to the initial four columns. To guarantee correct output tiles from WGMMA, the in-kernel transpose can similarly be tailored to emit a $\mathbf{V}$ tile featuring a correspondingly permuted set of rows.\footnote{Allowing for this flexibility in the in-kernel transpose obviates the need for shuffle instructions to transfer register values among threads, a method we discussed previously in~\citep{colfax_fp8_flashattention_2024}.}

\textbf{Accuracy: block quantization and incoherent processing.} In the FP8 (e4m3) format, 3 bits are dedicated to representing the mantissa, while 4 bits are reserved for the exponent. This limited bit allocation leads to increased numerical inaccuracies when compared to FP16 or BF16 formats. Furthermore, it is common for large-scale models to contain outlier values~\citep{dettmers2208llm, sun2024massive}—entries whose magnitudes greatly exceed the rest—posing additional challenges for quantization. The standard approach involves per-tensor scaling~\citep{micikevicius2022fp8}, wherein a single scaling factor is maintained for each tensor (such as a separate scalar for $\mathbf{Q}$, $\mathbf{K}$, and $\mathbf{V}$). To mitigate numerical imprecision arising in FP8 attention computations, we apply two strategies:

\begin{enumerate}

\item \textbf{Block quantization}: This method involves retaining a single scalar for each block, dividing the tensors $\mathbf{Q}$, $\mathbf{K}$, and $\mathbf{V}$ into blocks of dimension $B_r \times d$ or $B_c \times d$, and independently quantizing each block. The quantization step can be combined with any preceding operation in the attention pipeline, such as rotary embedding, without incurring additional computational latency (as rotary embedding is limited by memory bandwidth). Given that \textsc{FlashAttention-3} operates using blocks, the scaling for each block of $\mathbf{S}$ can be seamlessly adjusted for block quantization, introducing no extra computational burden.
\item \textbf{Incoherent processing}: To mitigate the impact of outliers, a random orthogonal matrix $\mathbf{M}$ is applied to $\mathbf{Q}$ and $\mathbf{K}$ prior to quantizing to FP8. Because orthogonality guarantees $\mathbf{M} \mathbf{M}^\top = I$, this means that $(\mathbf{Q}\mathbf{M})(\mathbf{K}\mathbf{M})^\top = \mathbf{Q}\mathbf{K}^\top$, so the ultimate attention computation is unaffected by this transformation. The effect of applying $\mathbf{M}$ is to randomize and disperse the values of $\mathbf{Q}$ and $\mathbf{K}$, as each element of $\mathbf{Q}\mathbf{M}$ or $\mathbf{K}\mathbf{M}$ becomes a random weighted sum of the original elements, thereby decreasing the overall quantization error. Consistent with \citet{chee2024quip} and~\citet{tseng2024quip}, the construction of $\mathbf{M}$ involves the product of random diagonal sign matrices ($\pm 1$) and a Hadamard matrix, facilitating efficient multiplication in $O(d \log d)$ time instead of $O(d^2)$ and enabling fusion with the rotary embedding process at negligible computational expense.
\end{enumerate}

Experimental results in \cref{sec:numerical_error} show that these methods can decrease numerical errors by up to a factor of 2.6.

\section{Empirical Validation}
\label{sec:exp}

Our implementation of \textsc{FlashAttention-3} leverages the WGMMA and TMA abstractions provided by CUTLASS~\citep{Thakkar_CUTLASS_2023} primitives, which we then assess in terms of both efficiency and accuracy.

\begin{itemize}

\item \textbf{Benchmarking attention.} We evaluate the performance of \textsc{FlashAttention-3} by recording its runtime across varying sequence lengths and benchmarking it against several alternatives: the conventional PyTorch implementation, \textsc{FlashAttention-2}, the Triton-based version of \textsc{FlashAttention-2} (which leverages H100-optimized instructions), and a proprietary cuDNN implementation of \textsc{FlashAttention-2} optimized for H100 GPUs. Our experiments show that \textsc{FlashAttention-3} delivers up to 2.0$\times$ speedup over \textsc{FlashAttention-2}, and is up to 1.5$\times$ faster compared to the Triton variant. Moreover, \textsc{FlashAttention-3} attains peak throughput of 740 TFLOPs/s, corresponding to 75\% of the H100 GPU's theoretical peak performance.
\item \textbf{Ablation study.} Our ablation analysis verifies that enhancements such as warp-specialization and GEMM-softmax pipelining are central to the performance gains achieved by \textsc{FlashAttention-3}.
\item \textbf{Accuracy of FP8 attention.} We demonstrate that the adoption of block quantization alongside incoherent processing yields a 2.6$\times$ reduction in the numerical error for FP8 \textsc{FlashAttention-3}.

\subsection{Benchmarking Attention}
\label{subsec:benchmark_attn}

We evaluate the performance of various attention mechanisms by measuring their runtimes on an H100 80GB SXM5 GPU under several configurations (with or without causal masking, and head dimensions of either 64 or 128), using FP16 inputs. The findings, detailed in~\cref{fig:benchmark_attn_fwd} and~\cref{fig:benchmark_attn_bwd}, indicate that \textsc{FlashAttention-3} achieves approximately 1.5-2.0$\times$ the speed of \textsc{FlashAttention-2} during the forward computation, and about 1.5-1.75$\times$ during the backward computation. In comparison with a conventional attention implementation, \textsc{FlashAttention-3} can achieve speedups in the range of 3-16$\times$. For moderate to long input sequences (at least 1k in length), \textsc{FlashAttention-3} even outperforms the specialized, proprietary cuDNN library, which is extensively optimized for H100 GPUs.

\paragraph{Benchmark settings:} The sequence length is adjusted across 512, 1k, up to 16k, while the batch size is chosen such that the aggregate token count remains at 16k. The hidden dimension is fixed at 2048, with head dimensions set to either 64, 128, or 256, corresponding to 32, 16, and 8 heads, respectively. For forward pass FLOPs computation, we employ:

\begin{equation*}
  4 \cdot \text{seqlen}^2 \cdot \text{head dimension} \cdot \text{number of heads}.
\end{equation*}

In the case of causal masking, we halve this value to reflect that nearly half of the computations are performed. To determine the FLOPs required for the backward pass, we take the forward pass FLOPs and multiply them by 2.5. This scaling arises from the forward pass involving 2 matrix multiplications, while the backward pass—owing to recomputation—requires 5 matrix multiplications.

Additionally, we evaluate the forward pass runtime for FP8 using comparable conditions. The findings for a headdim of 256 are presented in ~\cref{fig:benchmark_attn_fp8}, while comprehensive results can be found in \cref{sec:benchmark_attn_fp8_full}.

\subsection{Ablation Study: 2-Stage Pipelining Experiments}
\label{sec:2-stage-pipelining-experiments}

We conduct an ablation study on both the 2-stage WGMMA-softmax pipelining and the use of warp-specialization for the non-causal FP16 \textsc{FlashAttention-3}, maintaining fixed parameters $\{\text{batch}, \text{seqlen}, \text{nheads}, \text{hdim}\} = \{ 4, 8448, 16, 128\}$. The findings presented in~\cref{table:ablation_pipelining} demonstrate that the proposed algorithmic strategies—specifically, introducing asynchrony with warp-specialization and enabling overlap between GEMM and softmax—result in a notable increase in speed, achieving a performance boost from 570 TFLOPs to 661 TFLOPs.

\subsection{Numerical Error Validation}
\label{sec:numerical_error}

Given the growing concern around the numerical error associated with \textsc{FlashAttention}~\citep{golden2024flash}, we conduct a comparative analysis between \textsc{FlashAttention-2}, \textsc{FlashAttention-3}, and a conventional attention mechanism, using a high-precision FP64 reference as the benchmark. In order to model the presence of outlier features and activations observed in LLMs~\citep{dettmers2208llm, sun2024massive}, the elements of $\mathbf{Q}, \mathbf{K}, \mathbf{V}$ are sampled according to the distribution:

\begin{equation*}
  \mathcal{N}(0, 1) + \mathcal{N}(0, 100) \cdot \mathrm{Bernoulli}(0.001).
\end{equation*}

Specifically, all entries follow a normal distribution with mean zero and a standard deviation of 1, except for 0.1\% of them, to which we add an independent Gaussian noise with a standard deviation of 10. The root mean squared error (RMSE) outcomes are reported in~\cref{table:numerical_error}. For computations in FP16, both \textsc{FlashAttention-2} and \textsc{FlashAttention-3} deliver a 1.7$\times$ reduction in RMSE over the standard method, attributed to the use of FP32 precision for intermediate (softmax) results. The baseline FP8 attention approach employs per-tensor scaling, with accumulation during matrix multiplication in FP32 and intermediate softmax outputs stored in FP16. Leveraging block quantization and incoherent processing, \textsc{FlashAttention-3} in FP8 attains a 2.6$\times$ improvement in accuracy compared to this baseline.

\section{Dicussion, Limitations, Conclusion}
\label{sec:discussion}

Through the development of \textsc{FlashAttention-3}, we have shown that leveraging innovative programming strategies alongside hardware capabilities like asynchrony and reduced numerical precision can significantly enhance both the performance and reliability of attention mechanisms. Our method achieves a 1.5-2.0$\times$ improvement in attention speed relative to \textsc{FlashAttention-2}, and diminishes FP8 numerical error by 2.6$\times$ compared with conventional per-tensor quantization approaches. Nonetheless, certain aspects remain for future investigation: refining our approach for LLM inference, incorporating a persistent kernel architecture within the FP8 kernel,\footnote{In our evaluations, FP16 \textsc{FlashAttention-3} utilizes a persistent kernel along with a load balancing approach, whereas FP8 \textsc{FlashAttention-3} lacks these features. This difference partly accounts for the reduced performance of FP8 \textsc{FlashAttention-3} in the context of short sequence lengths and causal masking, as opposed to FP8 cuDNN kernels.} and examining the consequences of low-precision attention on extensive training workloads. Although our present study emphasizes Hopper GPUs, we anticipate that these methodologies will be broadly relevant to other hardware accelerators. Ultimately, we aspire that advancements in speed and fidelity of the attention primitive will facilitate progress in tasks involving longer contexts.

\end{document}